\documentclass[12pt]{ctexart}
\usepackage{a4}
\usepackage{graphicx}
\usepackage{geometry}
\usepackage{amsmath}
\usepackage{mathtools}
\usepackage{amssymb}
\geometry{a4paper,left=2cm,right=2cm,top=2.5cm,bottom=2.5cm}

\title{几何基础作业8}
\author{无名}
\date{\today}

\begin{document}
\maketitle
\section{在希尔伯特几何公理体系中叙述并证明外角定理}
\begin{figure}[htb]
    \begin{center}
        \includegraphics[width=0.5\linewidth]{01.png}
        \caption{}
    \end{center}
\end{figure}
\subsection{外角定理}
如图，在$\bigtriangleup ABC$中,延长$BC$至点$D$，那么$\angle CAB+\angle ABC=\angle DCA$
\subsection{证明}
根据II2(对于任意两点A和B，直线AB上至少有一点C，使得B在A、C之间)延长$BC$至$D$

作$CE \Vert  AB$，由IV平行公理(过定直线外一点，至多有一条直线与该直线平行)，那么$CE$唯一

由III4(一个角可以唯一地转移到另一条射线上，使其与给定角相等)，
假设$\angle ECD\neq \angle ABC$，作$\angle DCF=\angle DCE$FE在同一边，那么现在它们是同位角。

此时必须$CF\Vert AB$，但由IV平行公理(过定直线外一点，至多有一条直线与该直线平行)，这与$CE$唯一矛盾,故同位角相等，即$\angle ECD= \angle ABC$

同理，假设$\angle BAC \neq \angle ACE$，根据III4在E一侧作$\angle ACF=\angle BAC$,它们是内错角，故$AB\Vert CF$，再由IV得到矛盾，从而$\angle BAC= \angle ACE$

那么由于$\angle ACD=\angle ACE+\angle ECD=\angle ABC+\angle BAC$ 故外角定理成立

在三角形内不可能有角为零，故外角也一定大于任意一个不相邻内角。
\end{document}